% GENERATED by tools/build_m05_code_sheet.py -- do not edit.
\begin{codebox}{Box 1}{Build the karate club network}
import igraph as ig
import graph_tool.all as gt
karate = ig.Graph.Famous("Zachary")
g = gt.Graph(directed=False)
g.add_edge_list(karate.get_edgelist())
print(g.num_vertices(), g.num_edges())
\end{codebox}
\begin{linenotes}
\item[1] \texttt{import} loads a library: a bundle of ready-made tools. \texttt{igraph} is a library for networks. We use it for one thing only: Leiden. \texttt{as ig} is its short name.
\item[2] \texttt{graph\_tool} is the main library of this sheet. It fits stochastic block models and draws networks. \texttt{as gt} is its short name.
\item[3] \texttt{ig.Graph.Famous("Zachary")} builds a network that igraph already knows: Zachary's karate club. The \texttt{=} stores it under the name \texttt{karate}.
\item[4--5] graph-tool cannot read igraph's networks. It has its own kind, so we make an empty one (\texttt{directed=False} says that a friendship goes both ways) and fill it. \texttt{karate.get\_edgelist()} asks igraph for all the friendships as pairs of node numbers, and \texttt{add\_edge\_\allowbreak list} hands the pairs to graph-tool, which creates the nodes and the edges. This is the only bridge between the two libraries.
\item[6] \texttt{num\_vertices()} counts the nodes (the members) and \texttt{num\_edges()} counts the edges (the friendships). You should see \texttt{34 78}.
\end{linenotes}
